% Part: proof-theory
% Chapter: cut-elimination
% Section: midsequent

\documentclass[../../../include/open-logic-section]{subfiles}

\begin{document}

\olfileid{pt}{cut}{mid}

\olsection{મધ્ય-સિક્વન્ટ પ્રમેય અને હર્બ્રાન્ડનું પ્રમેય}

કટ-નિરાકરણ પ્રમેયનું મહત્ત્વનું
પરિણામ મધ્ય-સિક્વન્ટ પ્રમેય છે.
તે કહે છે કે જો
$\Gamma \Sequent \Delta$ની
એવી !!a{proof}~$\pi$ હોય
જેમાં દરેક !!{formula} પરિમાણક-અગ્ર
સ્વરૂપનું હોય, તો એવી
!!a{proof}~$\pi'$ છે જેમાં
$S \ident \Pi \Sequent \Lambda$
સિક્વન્ટ (\emph{મધ્ય-સિક્વન્ટ})
હોય અને~$\pi'$માં~$S$ની
ઉપરનાં બધાં અનુમાનો વિધાનાત્મક,
જ્યારે~$S$ની નીચેના બધાં અનુમાનો
પરિમાણકનાં હોય.
(!!^a{formula}~$!A$ ત્યારે
\emph{પરિમાણક-અગ્ર સ્વરૂપનું}
હોય જ્યારે તેનું સ્વરૂપ
\[\mathsf{Q}_1x_1\dots\mathsf{Q}_nx_n\,!B(x_1,\dots,x_n)\]
હોય, જ્યાં દરેક
$\mathsf{Q}_i$ એ~$\lforall$
અથવા~$\lexists$ છે.)
મધ્ય-સિક્વન્ટ પ્રમેયનું બીજું
મહત્ત્વનું પરિણામ હર્બ્રાન્ડનું
પ્રમેય છે. તેનો સામાન્ય કિસ્સો
વર્ણવવો થોડો મુશ્કેલ છે,
પણ વ્યાપક રીતે તે કહે છે કે
પ્રથમ-ક્રમ તર્કના દરેક
!!{provable} !!{sentence}~$!A$
માટે એવી વિધાનાત્મક
પુનરુક્તિ~$!A'$ છે જેમાંથી
$!A$ !!{prove}d કરી શકાય.
અહીં~$!B$ પરિમાણક-રહિત હોય ત્યારે
$\lexists[x][!B(x)]$ સ્વરૂપનાં
!!{sentence}s માટેનું રૂપ આપીએ:

\begin{thm}[હર્બ્રાન્ડનું પ્રમેય]
માનો કે~$!B(x)$ પરિમાણક-રહિત છે
અને $\Proves
\lexists[x][!B(x)]$.
તો એવાં પદો~$t_1$, \dots, $t_n$
છે કે
$\Proves !B(t_1) \lor \dots \lor !B(t_n)$.
\end{thm}

!!{formula}
$!B(t_1) \lor \dots \lor !B(t_n)$ને
$\lexists[x][!B(x)]$નો
\emph{હર્બ્રાન્ડ વિકલ્પ} કહે છે.
કારણ કે~$!B$ પરિમાણક-રહિત છે,
આ વિકલ્પ પણ પરિમાણક-રહિત છે.
આથી તે !!{provable} હોય તો
કટ-મુક્ત સાબિતીથી !!{provable} છે,
અને તેથી માત્ર વિધાનાત્મક
અનુમાનો વાપરતી !!a{proof}
પણ છે. એટલે તે પુનરુક્તિ છે.

હર્બ્રાન્ડના પ્રમેયનો મુશ્કેલ ભાગ
એ છે કે દરેક !!{provable}
!!{formula}~$\lexists[x][!B(x)]$
માટે હર્બ્રાન્ડ વિકલ્પ છે
તે બતાવવું. વ્યસ્ત દિશા,
એટલે~$\lexists[x][!B(x)]$નો
હર્બ્રાન્ડ વિકલ્પ હોય તો
તે સાબિત થઈ શકે, સહેલી છે.
ખરેખર હર્બ્રાન્ડ વિકલ્પમાંથી
$\lexists[x][!B(x)]$ ફલિત થાય છે.
દાખલા તરીકે~$n=2$ માટે
~$\Log{G3c}$માં આ
!!a{proof} છે:
\[
\Axiom$!B(t_1) \fCenter \lexists[x][!B(x)], !B(t_1), !B(t_2)$
\Axiom$!B(t_2) \fCenter \lexists[x][!B(x)], !B(t_1), !B(t_2)$
\RightLabel{\LeftR{\lor}}
\BinaryInf$!B(t_1) \lor !B(t_2) \fCenter \lexists[x][!B(x)], !B(t_1), !B(t_2)$
\RightLabel{\RightR{\lexists}}
\UnaryInf$!B(t_1) \lor !B(t_2) \fCenter \lexists[x][!B(x)], !B(t_2)$
\RightLabel{\RightR{\lexists}}
\UnaryInf$!B(t_1) \lor !B(t_2) \fCenter \lexists[x][!B(x)]$
\DisplayProof
\]

$\Sequent \lexists[x][!B(x)]$ની
!!a{proof}~$\pi$માંથી
હર્બ્રાન્ડ વિકલ્પ કાઢવા,
પહેલાં એવો કિસ્સો લો જેમાં
કટ-મુક્ત !!{proof}~$\pi$નાં
બધાં $\RightR{\lexists}$
અનુમાનો અંતે આવે છે:
\[
\AxiomC{}
\RightLabel{$\pi_1$}
\Deduce$ \fCenter \lexists[x][!B(x)], !B(t_1), \dots, !B(t_n)$
\RightLabel{\RightR{\lexists}}
\UnaryInf$ \fCenter \lexists[x][!B(x)], !B(t_2), \dots, !B(t_n)$
\RightLabel{$(n-1) \times \RightR{\lexists}$}
\Deduce$ \fCenter \lexists[x][!B(x)]$
\DisplayProof
\]
જો~$\pi$માં
$\lexists[x][!B(x)]$
સક્રિય હોય એવાં
$\RightR{\lexists}$
અનુમાનો આ બધાં જ હોય,
તો~$\pi_1$માં બીજું કોઈ
પરિમાણક-અનુમાન નથી.
\sourcecorrection{OLMD-221}{મૂળમાં અહીં અપરિભાષિત $\pi_1'$ લખાયું છે; ઉપરના વૃક્ષની ઉપરની ઉપસાબિતી $\pi_1$ છે.}
કારણ કે~$!B(x)$માં બીજાં
પરિમાણકો નથી અને અંતિમ
સિક્વન્ટના પૂર્વાંગમાં
કોઈ !!{formula}s નથી.
અર્થાત્,
$ \fCenter \lexists[x][!B(x)],
!B(t_1), \dots, !B(t_n)$
એ~$\pi$નો મધ્ય-સિક્વન્ટ છે.
વધુમાં, તેની ઉપર પરિમાણક-અનુમાનો
ન હોવાથી ઉપરના બધા સિક્વન્ટમાં
$\lexists[x][!B(x)]$ હોય છે,
પણ તે સક્રિય કે મુખ્ય નથી.
તેથી તેને~$\pi_1$ !!{proof}માંથી
કાઢી શકાય છે અને
$\fCenter !B(t_1), \dots, !B(t_n)$ની
!!a{proof} મળે છે.
પછી~$\RightR{\lor}$ના
$n-1$ ઉપયોગોથી
હર્બ્રાન્ડ વિકલ્પ
$\Sequent !B(t_1) \lor \dots \lor !B(t_n)$ની
!!a{proof} મળે છે.

મુશ્કેલી એ છે કે
$\Sequent \lexists[x][!B(x)]$ની
કટ-મુક્ત !!{proof}માં પણ
બધાં~\RightR{\lexists}
અનુમાનો એક જ અંતિમ જૂથમાં
હોય એમ ધારવું ચાલે નહીં.
તેમની વચ્ચે એવા અનુમાનો
હોઈ શકે જેમાં કોઈ~$!B(t_i)$
મુખ્ય હોય. તેથી મધ્ય-સિક્વન્ટ
પ્રમેય જોઈએ.

\begin{thm}[મધ્ય-સિક્વન્ટ પ્રમેય]\ollabel{thm:midsequent}
જો $\Gamma \Sequent \Delta$ની
\Log{G3c}-!!{proof}~$\pi$ હોય
અને~$\Gamma$ તથા~$\Delta$નાં
બધાં !!{formula}s પરિમાણક-અગ્ર
સ્વરૂપનાં હોય, તો એવી
!!a{proof}~$\pi'$ છે જેમાં
$\Pi \Sequent \Lambda$
સિક્વન્ટ છે અને તેની નીચેના
બધાં અનુમાનો પરિમાણકનાં,
ઉપરનાં બધાં વિધાનાત્મક છે.
\end{thm}

\begin{proof}
  ~$\pi$માંના પરિમાણક-અનુમાનનો
  \emph{ક્રમ} એટલે તેની નીચેના
  વિધાનાત્મક અનુમાનોની સંખ્યા,
  અને~$\pi$નો ક્રમ~$o(\pi)$ એટલે
  તેનાં બધાં પરિમાણક-અનુમાનોના
  ક્રમોનો સરવાળો, એમ વ્યાખ્યાયિત કરો.
  જો $o(\pi) = 0$, તો કોઈ
  પરિમાણક-અનુમાનની નીચે
  વિધાનાત્મક અનુમાન નથી,
  અને કામ પૂરું છે:
  પરિમાણક-અનુમાનો હોય તો
  બધાં એક-આધારવિધાનવાળાં હોવાથી
  સૌથી ઉપરનું એક પરિમાણક-અનુમાન
  છે; તેનું આધારવિધાન
  મધ્ય-સિક્વન્ટ છે.
  પરિમાણક-અનુમાન જ ન હોય તો
  અંતિમ સિક્વન્ટ મધ્ય-સિક્વન્ટ લઈએ.
  \sourcecorrection{OLMD-222}{મૂળ આધારકિસ્સો પરિમાણક-અનુમાન હોય ત્યારે જ સૌથી ઉપરના અનુમાનનું આધારવિધાન બતાવે છે; એકેય પરિમાણક-અનુમાન ન હોય તે કિસ્સામાં અંતિમ સિક્વન્ટ પૂરતું છે.}

  જો $o(\pi) > 0$, તો
  $>0$ ક્રમનું સૌથી નીચેનું
  પરિમાણક-અનુમાન પસંદ કરો.
  તેનો ક્રમ $>0$ હોવાથી
  તેની નીચે વિધાનાત્મક
  અનુમાન હોવું જોઈએ.
  તે સૌથી નીચેનું હોવાથી
  તેનો નિષ્કર્ષ બીજા
  પરિમાણક-અનુમાનનું આધારવિધાન
  હોઈ શકે નહીં; એવું હોય
  તો એ નીચેના પરિમાણક-અનુમાનની
  નીચે પણ વિધાનાત્મક અનુમાન
  હશે. કયું પરિમાણક-અનુમાન
  અને તેની નીચે કયું વિધાનાત્મક
  અનુમાન છે તેના (ઘણા)
  કિસ્સા જુદા પાડીએ.
  અંતિમ સિક્વન્ટમાં માત્ર
  પરિમાણક-અગ્ર સૂત્રો હોવાથી
  પરિમાણક-અનુમાનનું મુખ્ય
  સૂત્ર પછીના વિધાનાત્મક
  અનુમાનમાં સક્રિય હોઈ
  શકે નહીં. નહીં તો
  તે વિધાનાત્મક અનુમાનના
  મુખ્ય સૂત્રમાં વિધાનાત્મક
  સંયોજકના વ્યાપમાં
  પરિમાણક આવશે.
  ઉપ-!!{formula} ગુણધર્મથી
  તે અંતિમ સિક્વન્ટના કોઈ
  !!a{formula}નું
  ઉપ-!!{formula} હોવું પડશે.
  તેથી પરિમાણક-અનુમાનનું
  મુખ્ય !!{formula} નીચેના
  વિધાનાત્મક અનુમાનના સંદર્ભનું
  !!a{element} છે.

  બે ઉદાહરણો જોઈએ:

પહેલાં માનો કે $\RightR{\lforall}$નું
નિષ્કર્ષ $\LeftR{\lif}$ના જમણા
આધારવિધાનમાં આવે છે. ત્યારે~$\pi$ની
સંબંધિત ઉપ-!!{proof}નો અંત
ઉપ-!!{proof}~$\pi_1$માં થાય છે.
\[
\AxiomC{}
\RightLabel{$\pi_2$}
\Deduce$\Gamma' \fCenter \Delta', \lforall[x][!A(x)], !B$
\AxiomC{}
\RightLabel{$\pi_3$}
\Deduce$!C, \Gamma' \fCenter \Delta', !A(c)$
\RightLabel{\RightR{\lforall}}
\UnaryInf$!C, \Gamma' \fCenter \Delta', \lforall[x][!A(x)]$
\RightLabel{\LeftR{\lif}}
\BinaryInf$!B \lif !C, \Gamma' \fCenter \Delta', \lforall[x][!A(x)]$
\DisplayProof
\]
આપણે ધારીએ કે~$\pi$ નિયમિત છે;
આથી સ્વનિર્ધારક ચલ~$c$,
$\pi_3$ની બહાર આવતો નથી.
ખાસ કરીને તે $!B
\lif !C$માં કે~$\pi_2$માં
આવતો નથી. હવે સાબિતી~$\pi_1'$
જુઓ:
\[
\AxiomC{}
\RightLabel{$\pi_2'$}
\Deduce$\Gamma' \fCenter \Delta', \lforall[x][!A(x)], !A(c), !B$
\AxiomC{}
\RightLabel{$\pi_3$}
\Deduce$!C, \Gamma' \fCenter \Delta', \lforall[x][!A(x)], !A(c)$
\RightLabel{\LeftR{\lif}}
\BinaryInf$!B \lif !C, \Gamma' \fCenter \Delta', \lforall[x][!A(x)], !A(c)$
\RightLabel{\RightR{\lforall}}
\UnaryInf$!B \lif !C, \Gamma' \fCenter \Delta', \lforall[x][!A(x)], \lforall[x][!A(x)]$
\DisplayProof
\]
અહીં~$\pi_2$ના જમણા ભાગમાં
$!A(c)$ ઉમેરી શિથિલીકરણ
કરવાથી~$\pi_2'$ મળે છે.
(આથી કોઈ અનુમાન, ખાસ કરીને
ક્રમને અસર કરે એવું પરિમાણક-અનુમાન,
ઉમેરાતું નથી. વળી~$\pi_2$ની
સ્વનિર્ધારક ચલની કોઈ શરત પણ
ભંગ થતી નથી, કારણ કે~$!A(c)$માં
આવતો કોઈ અચળ~$\pi_2$માં
સ્વનિર્ધારક ચલ તરીકે વપરાતો નથી.)
$c$ એ~$!B \lif !C$માં
આવતો ન હોવાથી
\RightR{\lforall}ની
સ્વનિર્ધારક ચલની શરત
હજી પણ પૂરી થાય છે.

હવે સંકોચનના સ્વીકાર્યપણાથી
$!B \lif !C, \Gamma' \fCenter \Delta',
\lforall[x][!A(x)]$ની
!!a{proof}~$\pi_1''$ મળે છે.
$\lforall[x][!A(x)]$ માટે
$\RightR{\Contraction}$નું
સ્વીકાર્યપણું સાબિત કરતી
સાબિતી અને તેમાં વપરાયેલી
$\RightR{\lforall}$ની
ઉલટાવી શકાય તેવી હોવાની
સાબિતી કોઈ અનુમાનનો ક્રમ
બદલતી નથી: વધુમાં વધુ
અનુમાનો દૂર કરે છે.
આથી~$\pi_1''$માં
$\pi_1$ કરતાં વધુ
પરિમાણક-અનુમાનો નથી.
વળી~$\pi_1''$માં દરેક
પરિમાણક-અનુમાનની નીચે
તેનાં~$\pi_1$માંનાં
અનુરૂપ અનુમાન જેટલાં જ
વિધાનાત્મક અનુમાનો છે,
એક અપવાદ સિવાય:
છેલ્લા \RightR{\lforall}ની
નીચે~$\pi$માં \LeftR{\lif} હતું;
તેને~$\pi_1''$માં
ઉપર ખસેડ્યું છે.
હવે~$\pi$માં
$\pi_1$ને બદલે~$\pi_1''$
મૂકો. મળેલી !!{proof}નો
ક્રમ ઓછો છે, કારણ કે
ઓછામાં ઓછો
$\RightR{\lforall}$ અનુમાનનો
ક્રમ (ઓછામાં ઓછો)~$1$ ઘટ્યો છે.
હવે અનુગામી ધારણા લાગુ પડે છે.

પરિમાણક-અનુમાન
\RightR{\lexists} હોય તે
કિસ્સા વધુ સહેલા છે, કારણ કે
સંકોચન જરૂરી નથી અને
સ્વનિર્ધારક ચલની શરતની
ચિંતા પણ નથી. દાખલા તરીકે,
માનો કે~$\RightR{\lexists}$
પછી~$\RightR{\land}$ આવે છે
(આ વખતે તેના ડાબા
આધારવિધાન તરીકે).
સંબંધિત ઉપ-!!{proof}
$\pi_1$ આ છે:
\[
\AxiomC{}
\RightLabel{$\pi_2$}
\Deduce$\Gamma' \fCenter \Delta', \lexists[x][!A(x)], !A(t), !B$
\RightLabel{\RightR{\lexists}}
\UnaryInf$\Gamma' \fCenter \Delta', \lexists[x][!A(x)], !B$
\AxiomC{}
\RightLabel{$\pi_3$}
\Deduce$\Gamma' \fCenter \Delta', \lexists[x][!A(x)], !C$
\RightLabel{\RightR{\land}}
\BinaryInf$\Gamma' \fCenter \Delta', \lexists[x][!A(x)], !B \land !C$
\DisplayProof
\]
\sourcecorrection{OLMD-223}{મૂળ વૃક્ષમાં આ સિક્વન્ટના \(\Gamma'\) પહેલાં વધારાનું ઉદ્ગારચિહ્ન હતું; અનુમાનના આધારવિધાન અને નિષ્કર્ષ સાથે સુસંગત કરવા તે કાઢ્યું છે.}
હવે~$\pi$માં $\pi_1$ને
બદલે~$\pi_1'$ મૂકો:
\[
\AxiomC{}
\RightLabel{$\pi_2$}
\Deduce$\Gamma' \fCenter \Delta', \lexists[x][!A(x)], !A(t), !B$
\AxiomC{}
\RightLabel{$\pi_3'$}
\Deduce$\Gamma' \fCenter \Delta', \lexists[x][!A(x)], !A(t), !C$
\RightLabel{\RightR{\land}}
\BinaryInf$\Gamma' \fCenter \Delta', \lexists[x][!A(x)], !A(t), !B \land !C$
\RightLabel{\RightR{\lexists}}
\UnaryInf$\Gamma' \fCenter \Delta', \lexists[x][!A(x)], !B \land !C$
\DisplayProof
\]
અહીં~$\pi_3$ના જમણા ભાગમાં
$!A(t)$ ઉમેરી શિથિલીકરણ
કરવાથી~$\pi_3'$ મળે છે.
આ શિથિલીકરણમાં~$\pi_3$ના
દરેક સિક્વન્ટની જમણી બાજુ
માત્ર~$!A(t)$ ઉમેરાય છે,
તેથી કોઈ પરિમાણક-અનુમાનનો
ક્રમ બદલાતો નથી.
$\pi'$માં પરિમાણક-અનુમાનોના
ક્રમો~$\pi$ જેવા જ છે;
માત્ર~\RightR{\land} સાથે
અદલબદલ કરેલા
\RightR{\lexists} અનુમાનનો
ક્રમ~$1$ ઘટ્યો છે.
હવે અનુગામી ધારણા લાગુ પડે છે.
\end{proof}

\begin{prob}
\olref{thm:midsequent}ની
સાબિતીમાં એ કિસ્સા વર્ણવો
જ્યાં~\LeftR{\lexists}નું
નિષ્કર્ષ~\RightR{\lif}ના
આધારવિધાનમાં આવે છે,
અને જ્યાં~\LeftR{\lforall}નું
નિષ્કર્ષ~\LeftR{\lor}ના
ડાબા આધારવિધાનમાં આવે છે.
\end{prob}

\begin{proof}[હર્બ્રાન્ડના પ્રમેયની સાબિતી]
માનો કે~$\pi$નો અંત
$\Sequent \lexists[x][!A(x)]$માં
થાય છે.
\olref{thm:midsequent}થી
આપણે~$\pi$ને મધ્ય-સિક્વન્ટ~$S$
ધરાવતી !!a{proof}~$\pi'$માં
ફેરવી શકીએ છીએ. આ
મધ્ય-સિક્વન્ટનું સ્વરૂપ
\[\Sequent
\lexists[x][!A(x)], !A(t_1), \dots, !A(t_n).\]
હોવું જોઈએ.
$S$ની ઉપરનાં બધાં
અનુમાનો વિધાનાત્મક હોવાથી
$\lexists[x][!A(x)]$
$S$ની ઉપરના કોઈ અનુમાનમાં
મુખ્ય હોઈ શકે નહીં.
તે પરમાણ્વીય નથી, તેથી
સ્વયંસિદ્ધમાં પણ મુખ્ય
હોઈ શકે નહીં.
વળી~$\lexists[x][!A(x)]$,
$!A(t_1)$નું યોગ્ય
ઉપસૂત્ર હોઈ શકે નહીં
(યાદ રાખો કે~$!A(x)$
પરિમાણક-રહિત છે), એટલે
$S$ની ઉપર સક્રિય પણ
હોઈ શકે નહીં.
તેથી~$S$માં પૂરી થતી
ઉપ-!!{proof}માંથી
$\lexists[x][!A(x)]$
દૂર કરતાં
$\Sequent !A(t_1), \dots, !A(t_n)$ની
સાચી !!{proof} મળે છે.
તેમાંથી~$n-1$ વાર
$\RightR{\lor}$ લાગુ કરી
$\Sequent !A(t_1) \lor \dots \lor !A(t_n)$ની
!!a{proof} મેળવી શકાય છે.
\end{proof}

% \begin{prob}
% Find a Herbrand disjunction of $\lexists[x][\lforall[y][(!A(x) \lif
% !A(y))]]$ by finding a cut-free !!{proof} of $\Sequent
% \lexists[x][\lforall[y][(!A(x) \lif !A(y))]]$ and applying the
% transformation in \olref{thm:midsequent} to find a midsequent (if
% necessary).
% \end{prob}

\end{document}
